\BourbakiExercisesSection

\begin{enumerate}
\def\labelenumi{\arabic{enumi}.}
\item
  设 \(\top\) 是集合 E 上的一个合成律。用 F 表示 E 与由单个元素组成的集合 \(\{e\}\) 的和集（《集合论》II，§ 4，第 8 号），并把 E 与 \(\{e\}\) 和 F 中相应的子集等同起来，证明可以以一种且仅一种方式在 F 上定义一个合成律 \(\overline{\top}\) ，它在 E 上诱导出律 \(\top\) ，且在该律下 \(e\) 是单位元；若 \(\top\) 是结合的，则律 \(\overline{\top}\) 是结合的。（称 F 是由 E “通过添加一个单位元”而导出的。）
\item
  设 \(\mathsf{T}\) 是 E 上一个处处有定义的律。
\end{enumerate}

\begin{enumerate}
\def\labelenumi{(\alph{enumi})}
\item
  要使 \(\top\) 是结合的，必要且充分的是：每个左平移 \(\gamma_{x}\) 与每个右平移 \(\delta_{y}\) 在 E 到 E 的映射所成的集合中（配备律 \(f \circ g\) ）都可交换。
\item
  假设 E 有单位元。要使 T 是结合的且交换的，必要且充分的是：映射 \((x,y)\mapsto x\top y\) 是广群 \(E\times E\) 到广群 E 中的一个态射。
\end{enumerate}

\begin{enumerate}
\def\labelenumi{\arabic{enumi}.}
\setcounter{enumi}{2}
\item
  在 E 到 E 的映射所成的集合 F 中，关系 \(f \circ g = f \circ h\) 蕴含 g = h，必要且充分的是 f 是单射；关系 \(g \circ f = h \circ f\) 蕴含 g = h，必要且充分的是 f 是满射；f （在律 \(\circ\) 下）是可消去的，必要且充分的是 f 是双射。
\item
  要使集合E上存在一个合成法则，使得E的每一个置换都是在该法则下E到自身的同构，其充分必要条件是E含有0、1或3个元素。
\item
  对 \(2 \leqslant n \leqslant 5\) ，在一个具有 n 个元素的集合 E 上，确定所有处处有定义、具有单位元且使 E 的所有元素都可消去的律；对 n = 5，证明存在满足这些条件的非结合律。
\end{enumerate}

（注意：习题 6 至 17 涉及以乘法记法书写的结合律； \(e\) 表示单位元（当它存在时）， \(\gamma_a\) 与 \(\delta_a\) 表示由元素 \(a \in E\) 所作的平移；对每个子集 \(X\) （属于 \(E\) ），我们记 \(\gamma_a(X) = aX\) , \(\delta_a(X) = Xa\) 。）

\begin{enumerate}
\def\labelenumi{\arabic{enumi}.}
\setcounter{enumi}{5}
\item
  对于有限集合上的一个结合律，每个可消去元素都是可逆的（利用第 3 号，命题 3）。
\item
  给定集合 E 上的一个结合律与一个元素 \(x \in E\) ，设 A 为由 \(x^{n}\) （ \(n \in N^{*}\) ）组成的集合；若存在单位元，设 B 为由 \(x^{n}\) （ \(n \in N\) ）组成的集合；若进一步 x 可逆，设 C 为由 \(x^{a}\) （ \(a \in Z\) ）组成的集合。证明，若 A（相应地
\end{enumerate}

B，C）是无限的，则它（配备由 E 上给定的律所诱导的律）同构于带加法的 \(N^{*}\) （相应地，N，Z）。

\begin{enumerate}
\def\labelenumi{\arabic{enumi}.}
\setcounter{enumi}{7}
\item
  在习题7的记号中，设 A 是有限的；证明 A 包含一个且仅一个幂等元（§ 1，第 4 号）（注意，若 \(x^p\) 与 \(x^q\) 是幂等元，则 \(x^p = x^{pq}\) , \(x^q = x^{pq}\) ，因此 \(x^p = x^q\) ）；若 \(h = x^p\) ，则由诸 \(x^n\) （ \(n \geqslant p\) ）所成的集合是 E 的一个稳定子集 D，使得在 D 上诱导的律下，D 是一个群。
\item
  在集合 E 上的一个乘法律下，设 \(a\) 是 E 的一个左可消去的元素。
\end{enumerate}

\begin{enumerate}
\def\labelenumi{(\alph{enumi})}
\item
  若存在元素 \(u\) 使得 \(au = a\) ，证明 \(ux = x\) 对所有 \(x \in \mathbb{E}\) 成立；特别地，若 \(xu = x\) 对所有 \(x \in \mathbb{E}\) 成立，则 \(u\) 是单位元。
\item
  若存在 \(u \in E\) 使得 au = a，且存在 \(b \in E\) 使得 ab = u，证明 ba = u（考察 aba）；特别地，若存在单位元 e 及元素 b 使得 ab = e，则 b 是 a 的逆元。
\end{enumerate}

\begin{enumerate}
\def\labelenumi{\arabic{enumi}.}
\setcounter{enumi}{9}
\tightlist
\item
  若 a 与 b 是 E 的两个元素，且 ba 是左可消去的，证明 a 是左可消去的。由此推出，在 E 上的一个结合且交换的律下，E 的不可消去元素所成的集合 S 满足 ES ⊆ S（从而是稳定的）。
\end{enumerate}

¶ 11. 若 E 的每个元素 x 都是左可消去的，则称 E 为左半群。(a) 若 u 是幂等元（§ 1，第 4 号），则对所有 \(x \in E\) 有 ux = x（利用习题9(a)）；u 是 Eu 上诱导的律的单位元。

\begin{enumerate}
\def\labelenumi{(\alph{enumi})}
\setcounter{enumi}{1}
\item
  若 \(u\) 与 \(v\) 是 \(\mathbf{E}\) 的两个不同的幂等元，则 \(\mathrm{Eu} \cap \mathrm{Ev} = \varnothing\) ，且集合 \(\mathrm{Eu}\) 与 \(\mathrm{Ev}\) （配备由 \(\mathbf{E}\) 上的律诱导的律）是同构的。
\item
  设 R 为诸集合 Eu（u 遍历 E 的幂等元集合）的并集的补集。证明 ER ⊂ R；由此 R 是 E 的稳定子集。若 R 非空，则 \(aR \neq R\) 对所有 \(a \in R\) 成立（利用习题 9(a) 证明：否则 R 将包含一个幂等元）；特别地，此时 R 是无限的。R 称为左半群 E 的剩余。
\item
  若 R 非空且 E 中至少存在一个幂等元 u（即 \(E \neq R\) ），则对所有 \(x \in REu\) , \(xEu \neq Eu\) ；特别地，REu 中没有元素在 Eu 中可逆，且 REu 是无限集（利用习题 8）。
\item
  若 E 有单位元 e，则 e 是 E 中唯一的幂等元且 R 为空（注意 E = Ee）。
\item
  若 E 中存在一个右可消去的元素 \(a\) （特别地，若 E 上给定的律是交换的），则要么 E 有单位元，要么 \(\mathbf{E} = \mathbf{R}\) （注意，若存在幂等元 \(u\) ，则 \(xa = xua\) 对所有 \(x \in \mathbf{E}\) 成立）。例如，若 E 是整数 \(\geqslant 1\) 的集合并赋予加法，则 E 是一个交换半群且 \(\mathbf{E} = \mathbf{R}\) 。
\item
  若存在 \(a \in \mathbf{E}\) 使得 \(\delta_a\) 是满射的，则要么 \(\mathbf{E}\) 有单位元，要么 \(\mathbf{E} = \mathbf{R}\) （分别考察 \(a \in \mathbf{R}\) 的情形与 \(a \in \mathbf{Eu}\) （对某个幂等元 \(u\) ）的情形）。
\end{enumerate}

¶ 12. 在 E 上的一个乘法律下，设 \(a \in E\) 使得 \(\gamma_{a}\) 是满射的。

\begin{enumerate}
\def\labelenumi{(\alph{enumi})}
\item
  证明，若存在 \(u\) 使得 \(ua = a\) ，则 \(ux = x\) 对所有 \(x \in E\) 成立。
\item
  要使元素 \(b \in E\) 满足 ba 是左可消去的，必要且充分的是 \(\gamma_{a}\) 是满射的且 b 是左可消去的。
\end{enumerate}

¶ 13. 假设对所有 \(x \in E\) ， \(\gamma_x\) 都是满射的。证明若对某个元素 \(a \in E\) ， \(\gamma_a\) 是双射映射，则 \(\gamma_x\) 对所有 \(x \in E\) 都是双射映射（利用习题12(b)）；特别地，若存在两个元素 \(a, b\) （属于 \(E\) ）使得 \(ab = b\) （利用习题12(a)），则 \(E\) 是一个剩余 \(R\) 为空的左半群；此外，对每个幂等元 \(u\) ， \(Eu\) 的每个元素在 \(Eu\) 中可逆。

¶ 14. 对于有限集合 E 上的一个结合律，存在形如 aE 的极小子集（即按包含序排列的、具有这种形式的 E 的子集族中的极小元素）。

\begin{enumerate}
\def\labelenumi{(\alph{enumi})}
\item
  若 M = aE 是极小的，则对所有 \(x \in M\) 有 xM = xE = M；在由 E 上的律诱导的律下，M 是一个左半群（习题 11），其中每个左平移都是双射的（参见习题 13）。
\item
  若 \(\mathbf{M} = a\mathbf{E}\) 与 \(\mathbf{M}' = a'\mathbf{E}\) 是极小的且互不相同，则 \(\mathbf{M} \cap \mathbf{M}' = \varnothing\) ；对所有 \(b \in \mathbf{M}\) ，映射 \(x' \mapsto bx'\) （从 \(\mathbf{M}'\) 到 \(\mathbf{M}\) 中）是 \(\mathbf{M}'\) 到 \(\mathbf{M}\) 上的一个双射映射。由此推出，存在幂等元 \(u' \in \mathbf{M}'\) 使得 \(bu' = b\) ，以及幂等元 \(u \in \mathbf{M}\) 使得 \(uu' = u\) （取 \(u\) 为满足 \(bu = b\) 的幂等元）。证明 \(u'u = u'\) （考察 \(uu'u\) ），且每个 \(y' \in \mathbf{M}'\) 若满足 \(y'u = y'\) 便属于 \(\mathbf{M}'u'\) 。
\item
  证明映射 \(x' \mapsto ux'\) （从 \(M'u'\) 到 \(M\) 中）是 \(M'u'\) 到 \(Mu\) 上的一个同构；由此推出 \(M\) 与 \(M'\) 是同构的左半群。
\item
  设 \(M_{i}\) ( \(1 \leqslant i \leqslant r\) ) 是形如 aE 的各不相同的极小子集：由 (b) 推出，每个左半群 \(M_{i}\) 的幂等元可以排成一个序列 \((u_{ij})\) ( \(1 \leqslant j \leqslant s\) )，使得对所有 i、j、k 有 \(u_{ij}u_{kj} = u_{ij}\) 。若 K 是诸 \(M_{i}\) 的并集，证明 \(Eu_{ij} \subset K\) （注意，对所有 \(x \in E\) ， \(xu_{ij}E\) 是形如 aE 的极小子集）；由此推出 \(Eu_{ij}\) 是诸 \(u_{kj}Eu_{kj}\) （ \(1 \leqslant k \leqslant r\) ）的并集（证明 \((\mathrm{Eu}_{ij}) \cap \mathbf{M}_{k} = u_{kj}\mathrm{Eu}_{kj}\) ），且 \(Eu_{ij}E = K\) 。最后证明，形如 Eb 的每个极小子集都与这 s 个集合 \(Eu_{ij}\) 之一相同（注意，利用 (a)， \(Ebu_{ij}b = Eb\) ，并由此推出 \(Eb \subset K\) ）。
\end{enumerate}

\begin{enumerate}
\def\labelenumi{\arabic{enumi}.}
\setcounter{enumi}{14}
\tightlist
\item
  设 \((x_{i})_{1\leqslant i\leqslant n}\) 是左可消去元素的有限序列。
\end{enumerate}

\begin{enumerate}
\def\labelenumi{(\alph{enumi})}
\item
  证明关系 \(x_{1}x_{2}\ldots x_{n} = e\) 蕴含所有关系 \(x_{i + 1}\ldots x_nx_1x_2\ldots x_i = e\) （它们彼此由“循环置换”得到），对 \(1\leqslant i\leqslant n\) 。
\item
  由此推出，若序列 \((x_{i})\) 的复合是可逆的，则每个 \(x_{i}\) 都是可逆的。
\end{enumerate}

¶ 16. 设 \(\top\) 是集合 E 上的一个结合律，E* 为 E 的可消去元素所成的集合；假设 E* 非空且每个可消去元素都是中心元素。设 \(\mathfrak{F}\) 表示具有以下性质的子集 \(\mathbf{X}\) （ \(\mathbf{E}\) 的）所成的集合：存在 \(y \in \mathrm{E}^*\) 使得 \(\delta_y(\mathrm{E}) \subset \mathbf{X}\) 。

\begin{enumerate}
\def\labelenumi{(\alph{enumi})}
\item
  证明 \(\mathfrak{F}\) 中两个集合的交集属于 \(\mathfrak{F}\) 。
\item
  设 \(\Phi\) 为这样的函数 \(f\) 所成的集合：定义在 \(\mathfrak{F}\) 中一个集合上，取值于
\end{enumerate}

E 且满足：对 \(X \in \mathfrak{F}\) ， \(f(\mathbf{X})\) 属于 F。设 R 表示 \(\Phi\) 的元素 f 与 g 之间的关系：“存在集合 \(X \in F\) 使得 f 与 g 在 X 上的限制相同”。证明 R 是一个等价关系；设 \(\Psi = \Phi/R\) 为 \(\Phi\) 在该关系下的商集。

\begin{enumerate}
\def\labelenumi{(\alph{enumi})}
\setcounter{enumi}{2}
\item
  设 \(f\) 与 \(g\) 是 \(\Phi\) 的两个元素， \(A \in \mathfrak{F}\) 与 \(B \in \mathfrak{F}\) 为 \(f\) 与 \(g\) 分别定义于其上的集合；存在 \(X \subset B\) 属于 \(\mathfrak{F}\) ，使得 \(g(X) \subset A\) ；若 \(g_X\) 是 \(g\) 在 \(X\) 上的限制，证明映射 \(f \circ g_X\) 属于 \(\Phi\) ，且其类（模 R）仅依赖于 \(f\) 与 \(g\) 的类（而不依赖于 \(X\) ）；这个类称为 \(f\) 的类与 \(g\) 的类的复合；证明这样在 \(\Psi\) 上定义的合成律是结合的且有单位元。
\item
  对所有 \(a \in \mathbb{E}\) ，证明左平移 \(\gamma_{a}\) 属于 \(\Phi\) ；令 \(\phi_{a}\) 为其类（模 R）。证明映射 \(x \mapsto \phi_{x}\) 是 E 到 \(\Psi\) 的一个子幺半群上的同构，且若 \(x \in \mathbb{E}^{*}\) ，则 \(\phi_{x}\) 在 \(\Psi\) 中可逆（考察 \(\gamma_{x}\) 的逆映射，证明它属于 \(\Phi\) 且其类（模 R）是 \(\phi_{x}\) 的逆元）。由此推出第 4 号定理 1 的一个推广。
\end{enumerate}

\begin{enumerate}
\def\labelenumi{\arabic{enumi}.}
\setcounter{enumi}{16}
\tightlist
\item
  设E是一个幺半群，S是E的一个稳定子集，满足以下性质：
\end{enumerate}

\((\alpha)\) 对所有 \(s \in S\) 以及所有 \(a \in E\) ，存在 \(t \in S\) 与 \(b \in E\) 使得 \(sb = at\) 。

(β) 对所有 \(a, b \in \mathbb{E}\) 以及所有 \(s \in \mathbb{S}\) （满足 \(sa = sb\) ），存在 \(t \in \mathbb{S}\) 使得 \(at = bt\) 。

\begin{enumerate}
\def\labelenumi{(\alph{enumi})}
\tightlist
\item
  在 E × S 中用 \((a, s) \sim (b, t)\) 记关系：“存在 S 中的 \(s'\) 与 \(t'\) 使得 \(tt' = ss'\) 且 \(as' = bt'\) 。”
\end{enumerate}

证明 \(\sim\) 是一个等价关系。我们记 \(\overline{\mathbf{E}} = \mathbf{E} \times \mathbf{S}/\sim\) ，用 \(as^{-1}\) 表示 \((a, s) (\text{mod } \sim)\) 的等价类，并用 \(\varepsilon: \mathbf{E} \to \overline{\mathbf{E}}\) 表示映射 \(a \mapsto ae^{-1}\) 。

\begin{enumerate}
\def\labelenumi{(\alph{enumi})}
\setcounter{enumi}{1}
\item
  证明在 \(\overline{\mathbf{E}}\) 上存在一种且仅一种幺半群结构，使得 \(\varepsilon\) 是保单位同态，且对所有 \(s \in \mathbf{S}\) ， \(\varepsilon(s)\) 是可逆的。
\item
  证明 \((\overline{\mathbf{E}},\varepsilon)\) 具有定理 1（第 4 号）中所述的泛性质。
\item
  证明 \(\varepsilon\) 是单射当且仅当S中的元素都是可消去的。
\end{enumerate}

